\section{Multimodal integration design}

To evaluate multimodal representations and fusion strategies consistently, we adopt a unified integration architecture across the evaluated models, providing a common interface that maps heterogeneous visual and textual features into a shared latent space. We distinguish two integration paradigms according to where modality fusion occurs: data-level integration and model-level integration.

Data-level integration, exemplified by architectures such as BeLightRec \parencite{10.1007/978-981-97-9616-8_3}, performs multimodal fusion offline at the data-loading layer. The raw visual features $\mathbf{x}_i^{\text{vis}}$ (from MobileNetV2 or CLIP) and textual features $\mathbf{x}_i^{\text{txt}}$ (from TF-IDF) are combined to compute an item-item similarity matrix---for late fusion (\texttt{multimodal}), over concatenated normalized representations---which is then filtered, normalized, and cached. The model does not ingest raw modal embeddings during training; instead, the precomputed similarity matrix $S$ acts as a static structural prior guiding graph propagation. We couple this design with the dual-graph propagation of CombiGCN, decoupling graph convolutions from high-dimensional feature processing and thereby reducing the GPU memory and compute footprint during training.

Conversely, model-level integration is employed by BM3 \parencite{10.1145/3543507.3583251} and FREEDOM \parencite{Zhou_2023}. In this paradigm, raw multimodal embeddings ($\mathbf{x}_i^{\text{vis}} \in \mathbb{R}^{d_v}$ and $\mathbf{x}_i^{\text{txt}} \in \mathbb{R}^{d_t}$) are loaded directly into the GPU memory, and the fusion process is learned end-to-end. To address the discrepancy in dimensions between visual features (e.g., $d_v = 512$ for CLIP and $d_v = 768$ for MobileNetV2 after PCA reduction) and textual features (e.g., $d_t = V$ for TF-IDF, where $V$ is the vocabulary size), we implement linear projection layers mapping features to a shared latent space:
\begin{equation}
    \mathbf{h}_{v,i} = \mathbf{W}_v \mathbf{x}_i^{\text{vis}} + \mathbf{b}_v, \quad \mathbf{h}_{t,i} = \mathbf{W}_t \mathbf{x}_i^{\text{txt}} + \mathbf{b}_t,
\end{equation}
where $\mathbf{W}_v \in \mathbb{R}^{d \times d_v}$ and $\mathbf{W}_t \in \mathbb{R}^{d \times d_t}$ are learnable projection matrices, $\mathbf{b}_v, \mathbf{b}_t \in \mathbb{R}^d$ are bias vectors, and $d = 512$ is the target dimension matching the user-item collaborative filtering embedding space. 

After projection, the features are integrated depending on the configured fusion strategy. For the late fusion (\texttt{multimodal}) strategy, the representation is computed via element-wise averaging:
\begin{equation}
    \mathbf{e}_{i, \text{modal}} = \frac{\mathbf{h}_{v,i} + \mathbf{h}_{t,i}}{2}.
\end{equation}
For the weighted attention fusion (\texttt{multimodal\_attention}) strategy, we implement an adaptive attention gating network. The projected visual and textual embeddings are concatenated and passed through a trainable attention layer to compute dynamic modality weights:
\begin{equation}
    \beta_{v,i} = \sigma\left(\mathbf{w}_{\text{attn}}^{T} [\mathbf{h}_{v,i} \parallel \mathbf{h}_{t,i}] + b_{\text{attn}}\right),
\end{equation}
\begin{equation}
    \beta_{t,i} = 1 - \beta_{v,i},
\end{equation}
where $\parallel$ denotes the vector concatenation operator, $\mathbf{w}_{\text{attn}} \in \mathbb{R}^{2d}$ is a trainable weight vector, $b_{\text{attn}} \in \mathbb{R}$ is a bias scalar, and $\sigma(\cdot)$ is the sigmoid activation function. The fused multimodal representation is then defined as:
\begin{equation}
    \mathbf{e}_{i, \text{modal}} = \beta_{v,i} \mathbf{h}_{v,i} + \beta_{t,i} \mathbf{h}_{t,i}.
\end{equation}
This gating mechanism enables the recommendation models to dynamically prioritize the most informative modality for each individual fashion item, balancing visual aesthetics and textual descriptions adaptively. The resulting $\mathbf{e}_{i, \text{modal}}$ is then integrated into the GNN collaborative filtering representations via contrastive learning objectives (bootstrap alignment in BM3 or InfoNCE in FREEDOM) and combined directly into the final preference scoring function.
